\documentclass[12pt,a4paper]{article}
\usepackage[utf8]{inputenc}
\usepackage[T1]{fontenc}
\usepackage[english]{babel}
\usepackage[a4paper, margin=2.5cm]{geometry}
\usepackage{hyperref}
\usepackage{parskip}
\pagestyle{empty}

\setlength{\parskip}{4pt}
\setlength{\parindent}{0pt}

\begin{document}

\begin{center}
\textbf{\Large IA382 Final Project Proposal}

\vspace{4pt}
{\small Felipe Rezende Gardin RA: 215766 \quad \today}
\end{center}

\vspace{12pt}

\noindent\textbf{Project Title}

Automated Presentation Generation: Transforming Seminar Transcripts into Professional Slide Decks Using AI

\vspace{8pt}

\noindent\textbf{Project Abstract}

This project leverages multiple AI tools to automatically transform academic seminar transcripts into polished
PowerPoint presentations. The methodology involves three stages: (1) using Claude API to parse the transcript and
extract key concepts, narrative structure, and content hierarchy; (2) generating a structured outline that feeds into
SlidesAI.io for automated slide design and layout; (3) refining the output through manual curation of visuals and speaker notes.
As a concrete case study, this project will demonstrate the pipeline using the ``Human-AI Collaboration'' seminar
transcript by Neeli Prasad from this IA382 course, transforming the raw lecture content into a professional 20--30 slide
presentation ready for educational use. This project demonstrates how orchestrating multiple AI tools can dramatically reduce
the time required to transform raw academic content into structured, visually compelling knowledge artifacts.

\vspace{8pt}

\noindent\textbf{Target Seminar}

Human-AI Collaboration by Neeli Prasad (IA382-FEEC-UNICAMP Seminar Series). This seminar explores ethical frameworks, 
privacy concerns, and best practices for human-AI collaboration in real-world applications---topics directly relevant 
to the responsible orchestration of multiple AI tools proposed in this project.

\vspace{8pt}

\noindent\textbf{AI Tools to Be Used}

\begin{itemize}
  \item Claude API: transcript analysis, concept extraction, and content structuring
  \item SlidesAI.io: automated slide design and layout generation
  \item Notebook LM (optional): knowledge synthesis and deeper conceptual mapping
  \item Perplexity (optional): research and fact verification for accuracy validation
\end{itemize}

\end{document}
